%% put this in main.tex and it will work

\documentclass{article}

% load your entire custom package library with one line.
\usepackage{bsk_latex_lib/bsk_latex_lib}

% --- document configuration ---
% all high-level settings for this specific document are placed here.
\setupprotokoll{
  sektion = {root},
  plats   = {karlskrona/microsoft teams},
  datum   = {2025-mm-dd},
}

% --- person database ---
% define all relevant people and their roles.
\newperson{\ordfarande} { name = {ordförande},    role = {ordförande} }
\newperson{\person}  { name = {namn namnson},     role = {vice-ordförande} }
\newperson{\namn} { name = {namn namnson},       role = {sekreterare} }
% \newperson{\namn} { name = {namn namnson}, role = {ekonomiansvarig} }
% \newperson{\namn} { name = {namn namnson},   role = {ledamot} }
% \newperson{\namn}  { name = {namn namnson},   role = {ledamot} }
% \newperson{\namn} { name = {namn namnson}, role = {ledamot} }
% \newperson{\namn} { name = {namn namnson},    role = {suppleant} }
% \newperson{\namn} { name = {namn namnson},   role = {suppleant} }

\begin{document}

% --- meeting-specific data ---
% this is the only section that needs to be filled in for each new meeting.

% 1. clear any data from previous meetings (good practice).
\clearattendance

% 2. assign the officials for this meeting.
\setmeetingofficials{\ordfarande}{\person}{\namn} % chairperson, secretary, adjuster

% 3. set the attendance for this meeting.
\addpresent{\ordfarande}
\addpresent{\person}
\addpresent{\namn}

% \addabsent{\namn}
% \addabsent{\namn}

% --- automated document generation ---

% this single command now generates the attendance lists and all
\printformalia

% --- meeting content (ärenden) ---
% the unique content of the meeting is written here.

\paragraf{skrivelser}{
  mail från xx angående yy \\
  mail från xx angående yy \\
}

\paragraf{rapporter}{
  % updated: use \person{\melker} to get the name
  \person{\ordfarande} har:
  \begin{itemize}
    \item inget att rapportera.
  \end{itemize}
  \person{\namn} har:
  \begin{itemize}
    \item inget att rapportera.
  \end{itemize}
}

\paragraf{tidigare action-points}{}

\paragraf{studiebevakning}{
  \textbf{möte som vi kanske varit på}\\
  brödtext
  \enter

  \textbf{annat möte eller punkt}\\
  brödtext
}
\paragraf{feedback-soffan}{
  inget att rapportera.
}

\paragraf{studiesocialt}{
  \textbf{event}\\
  brödtext
  \enter

  \textbf{annat event}\\
  brödtext
}

\paragraf{styrelsen}{}
\paragraf{vibe check}{}

% --- closing section ---
% the \printclosing command handles the decision summary and signatures.
% the paragraphs for övrigt, kalender, etc., are placed before it.

\paragraf{övrigt}{}
\paragraf{kalender}{
  \begin{tabular}{ll}
  nästa styrelse möte & yyyy-mm-dd\\
  event 1 & yyyy-mm-dd\\
  event 2 & yyyy-mm-dd\\
  möte 1 & yyyy-mm-dd\\
  möte 2 & yyyy-mm-dd\\
  \end{tabular}
}

% §15 for beslutssammanfattning is now generated automatically by \printclosing.

\paragraf{action points}{
  \person{\ordfarande} ska:
  \begin{itemize}
    \item göra x
    \item göra y
  \end{itemize}
  \person{\namn} ska:
  \begin{itemize}
    \item maila x
    \item göra klart y
  \end{itemize}
}


\printclosing

\end{document}
